\chapter{Bilan, validation et perspectives}

\section{Stratégie de validation et indicateurs de suivi}

La validation d'une évolution d'alerting ne se limite pas à la réussite technique du déploiement : elle
doit mesurer la pertinence opérationnelle du système obtenu. J'ai donc défini un ensemble
d'indicateurs à suivre dans la durée :

\begin{itemize}
	\item le volume d'alertes par niveau de criticité ;
	\item le ratio entre alertes actionnables et alertes non pertinentes ;
	\item le temps moyen de détection (MTTD) et de réparation (MTTR) ;
	\item le ressenti des équipes d'astreinte, en particulier leur niveau de fatigue et leur confiance dans les alertes.
\end{itemize}

Ces métriques permettent d'ajuster continûment les seuils et la classification, dans une logique
d'amélioration continue.

\section{Gains attendus de la solution}

La solution conçue apporte des bénéfices à plusieurs niveaux :

\begin{itemize}
	\item \textbf{Gestion des astreintes et réactivité} : réduction de la fatigue d'alerte — seules les alertes urgentes ou critiques réveillent la personne d'astreinte —, amélioration des temps de détection et de réparation en identifiant directement le service problématique, et fin des interventions inutiles.
	\item \textbf{Infrastructure et développement} : visibilité précise sur la plateforme, aidant les développeurs à savoir exactement quels services nécessitent une intervention.
	\item \textbf{Expérience utilisateur et entreprise} : amélioration de la disponibilité de l'application grâce à une résolution plus rapide des incidents, ce qui contribue à la réputation de Winamax, à la diminution des plaintes utilisateurs et donc à son attractivité.
\end{itemize}

\section{Limites}

Le travail réalisé constitue une base méthodologique et technique solide, mais plusieurs limites
doivent être soulignées. La liste des métriques d'application, aujourd'hui centrée sur trois indicateurs
clés, est vouée à évoluer et à s'affiner service par service. La gestion des alertes personnalisées n'est
pas encore prise en charge par le module Terraform. Le calcul de l'Error Rate par méthode, envisagé,
reste délicat car les seuils diffèrent selon les méthodes et pourraient réintroduire du bruit. Enfin, la
pleine réussite de la solution dépend d'une gouvernance rigoureuse : sans réévaluation régulière des
seuils et sans discipline collective, un système d'alerting tend naturellement à se dégrader et à
redevenir bruyant.

\section{Perspectives}

Plusieurs axes d'évolution se dégagent :

\begin{itemize}
	\item affiner la granularité des niveaux de criticité selon les services et étendre la surveillance aux composants backend (workers, Kafka, Redis), avec par exemple des alertes sur le consumer lag ;
	\item renforcer l'intégration avec les runbooks et les dashboards Grafana, pour un diagnostic toujours plus rapide ;
	\item automatiser le cycle de vie des alertes via Backstage, afin de garantir une cohérence durable entre les services déployés et leur configuration d'alerting ;
	\item industrialiser le suivi de la qualité des alertes — tableaux de bord d'indicateurs, revues régulières — afin de pérenniser les gains obtenus.
\end{itemize}

\section{Bilan personnel}

Sur le plan personnel, ce stage a été une expérience extrêmement enrichissante. Il m'a permis de
découvrir de l'intérieur le fonctionnement d'une équipe Infrastructure au sein d'une entreprise à forte
culture technique, et de contribuer à un projet à réel impact opérationnel. J'ai consolidé mes
compétences en cloud AWS, en Infrastructure as Code (Terraform), en observabilité (Prometheus,
CloudWatch, Grafana) et en ingénierie de fiabilité (méthodes RED et USE, niveaux de criticité, notions
de MTTD et MTTR). Au-delà de la technique, j'ai appris à conduire un projet dans un environnement de
production exigeant, à documenter et défendre mes choix de conception, et à collaborer avec des
équipes aux besoins variés. Cette expérience a confirmé mon intérêt profond pour le domaine de
l'infrastructure et du DevOps, et m'a donné une vision concrète du métier vers lequel je souhaite
m'orienter.

\section{Auto-positionnement : compétences mobilisées et projet professionnel}

Le tableau \ref{tab:synthese-competences} synthétise les compétences du référentiel R\&T mobilisées
pendant ce stage.

\begin{table}[h]
	\centering
	\caption{Synthèse des compétences exercées durant le stage}
	\label{tab:synthese-competences}
	\begin{tabular}{|p{0.3\textwidth}|p{0.6\textwidth}|}
		\hline
		\textbf{Compétence du référentiel} & \textbf{Mise en œuvre durant le stage} \\
		\hline
		Planification & Démarche itérative (immersion, spécification, développement, transition) ; phasage du déploiement en trois étapes. \\
		\hline
		Développement & Module Terraform pour la gestion des alertes CloudWatch/AMP ; règles Prometheus (promrules). \\
		\hline
		Documentation & Spécification technique du modèle RED/USE, runbooks par service, rédaction de ce rapport. \\
		\hline
		Évaluation en mise en œuvre & Définition d'indicateurs de suivi (volume d'alertes, ratio actionnable, MTTD/MTTR) et de la stratégie de double-run. \\
		\hline
		Administration \& exploitation & Branchement à incident.io (grouping, routing, escalade), gestion du cycle de vie des alertes ; missions de fast (provisionnement, réseau, Puppet). \\
		\hline
		Gestion des dysfonctionnements & Diagnostic de bout en bout des incidents réseau et des pannes de la chaîne d'alerting durant le fast. \\
		\hline
	\end{tabular}
\end{table}

Ce stage confirme mon projet professionnel : m'orienter vers l'ingénierie infrastructure et le DevOps.
Les compétences développées en IaC, en observabilité et en conduite de projet en production
constituent un socle réutilisable pour la suite de mon parcours d'ingénieur R\&T.
